% Preambulo compartido para clases en formato beamer
\documentclass[9pt, xcolor=svgnames]{beamer}
\usepackage[utf8]{inputenc}
\usepackage[T1]{fontenc}
\usepackage[spanish,es-nodecimaldot]{babel}

\usetheme[progressbar=frametitle]{metropolis}
\metroset{sectionpage=none, subsectionpage=none}
\usepackage{appendixnumberbeamer}
\usepackage{booktabs}
\usepackage{amsmath,amssymb,bm,mathtools}
\usepackage{physics}
\usepackage{graphicx}
\usepackage{multicol}
\usepackage{tikz}
\usepackage{minted}
\usemintedstyle{tango}
\setminted{fontsize=\scriptsize, breaklines=true, linenos=true, bgcolor=LightGray!20, frame=lines}
\newminted{python}{fontsize=\scriptsize, breaklines=true, linenos=true, bgcolor=LightGray!20, frame=lines}
\usepackage{pgfplots}
\pgfplotsset{compat=1.18}
\usepackage{ragged2e}
\usepackage{xspace}
\newcommand{\themename}{\textbf{\textsc{metropolis}}\xspace}

\setbeamercolor{block title}{bg=MidnightBlue!85, fg=white}
\setbeamercolor{block body}{bg=MidnightBlue!5, fg=black}
\setbeamercolor{alerted text}{fg=Crimson}
\setbeamercolor{frametitle}{bg=MidnightBlue!10, fg=black}
\setbeamertemplate{blocks}[rounded][shadow=true]


\weektitle{08}{Reaccion-difusion y formacion de patrones}

\begin{document}

\frame{\titlepage}

\begin{frame}{Organizacion de la semana}
\begin{block}{Formato unificado}
La Semana 08 se divide en \textbf{Clase 1} y \textbf{Clase 2}, manteniendo la estructura habitual del curso:
\begin{itemize}
    \item \textbf{30--45 min} de teoria y discusion conceptual.
    \item \textbf{75--90 min} de laboratorio guiado con Python.
\end{itemize}
\end{block}

\vspace{0.2cm}
\begin{columns}[T,onlytextwidth]
\column{0.48\textwidth}
\textbf{Clase 1}
\begin{itemize}
    \item ecuacion de difusion y transporte;
    \item discretizacion espacio-temporal;
    \item condicion de estabilidad;
    \item simulacion de difusion 1D y 2D.
\end{itemize}

\column{0.48\textwidth}
\textbf{Clase 2}
\begin{itemize}
    \item sistemas reaccion-difusion;
    \item inestabilidad de Turing;
    \item modelo de Gray--Scott;
    \item patrones, metricas y exploracion de parametros.
\end{itemize}
\end{columns}
\end{frame}

\section{Clase 1}

\begin{frame}[plain]
\centering
\vfill
{\Huge \bfseries Clase 1}\\[0.4cm]
{\Large Difusion, discretizacion y simulacion numerica}
\vfill
\end{frame}

\begin{frame}{Objetivos de la Clase 1}
\begin{itemize}
    \item Entender la difusion como un proceso de suavizado espacial impulsado por gradientes.
    \item Conectar una ecuacion diferencial parcial con una regla computacional local.
    \item Implementar un esquema explicito en una dimension y extenderlo a dos dimensiones.
    \item Discutir el rol del paso temporal y de la malla espacial en la estabilidad numerica.
\end{itemize}
\end{frame}

\begin{frame}{Motivacion fisica}
\justifying
La difusion aparece cuando una cantidad se redistribuye desde regiones donde esta muy concentrada hacia regiones donde esta menos concentrada.

\begin{itemize}
    \item calor en una barra;
    \item tinta o contaminantes en un fluido;
    \item densidad superficial de particulas;
    \item mezcla y relajacion hacia estados mas uniformes.
\end{itemize}

\begin{alertblock}{Idea clave}
La dinamica local depende del entorno cercano: una region con exceso de concentracion tiende a perder material hacia sus vecinos.
\end{alertblock}
\end{frame}

\begin{frame}{Ley de Fick y ecuacion de difusion}
\justifying
Si $u(x,t)$ representa concentracion, temperatura o densidad, la ley de Fick propone que el flujo es proporcional al gradiente:
\[
\bm{J} = -D \nabla u,
\]
donde $D$ es el coeficiente de difusion.

Combinando esta ley con conservacion obtenemos
\[
\pdv{u}{t} = D \nabla^2 u.
\]

\begin{block}{Interpretacion}
El laplaciano compara el valor local con sus alrededores. Si un punto esta muy por encima del promedio de sus vecinos, tendera a disminuir.
\end{block}
\end{frame}

\begin{frame}{Caso unidimensional}
En una dimension, la ecuacion toma la forma
\[
\pdv{u}{t} = D \pdv[2]{u}{x}.
\]

\begin{columns}[T,onlytextwidth]
\column{0.57\textwidth}
\begin{itemize}
    \item si $u(x,t)$ es un pulso angosto, con el tiempo se ensancha;
    \item la altura maxima disminuye;
    \item la masa total puede conservarse si no hay flujo por los bordes.
\end{itemize}

\begin{block}{Pregunta computacional}
\textquestiondown Como traducimos una derivada espacial a operaciones entre elementos de un arreglo?
\end{block}

\column{0.38\textwidth}
\begin{tikzpicture}[scale=0.9]
\begin{axis}[
    width=5.0cm,
    height=4.2cm,
    xmin=-3,xmax=3,
    ymin=0,ymax=1.2,
    axis lines=left,
    xtick=\empty,
    ytick=\empty
]
\addplot[very thick,MidnightBlue,domain=-3:3,samples=150]{exp(-4*x^2)};
\addplot[very thick,Crimson,domain=-3:3,samples=150]{0.7*exp(-x^2)};
\end{axis}
\end{tikzpicture}
\end{columns}
\end{frame}

\begin{frame}{Discretizacion por diferencias finitas}
\justifying
Tomamos una grilla uniforme:
\[
x_i = i\,\Delta x, \qquad t_n = n\,\Delta t.
\]

Usamos las aproximaciones
\[
\pdv{u}{t} \approx \frac{u_i^{n+1}-u_i^n}{\Delta t},
\qquad
\pdv[2]{u}{x} \approx \frac{u_{i+1}^n - 2u_i^n + u_{i-1}^n}{(\Delta x)^2}.
\]

Entonces el esquema explicito queda
\[
u_i^{n+1} = u_i^n + r\left(u_{i+1}^n - 2u_i^n + u_{i-1}^n\right),
\qquad
r = \frac{D\,\Delta t}{(\Delta x)^2}.
\]
\end{frame}

\begin{frame}{Lectura fisica del esquema}
\begin{columns}[T,onlytextwidth]
\column{0.58\textwidth}
\begin{itemize}
    \item Si $u_i^n$ es mayor que el promedio de sus vecinos, el termino entre parentesis es negativo.
    \item Si $u_i^n$ es menor que el promedio de sus vecinos, ese termino es positivo.
    \item El algoritmo tiende a suavizar gradientes fuertes.
\end{itemize}

\begin{alertblock}{Mensaje}
La discretizacion no solo aproxima una ecuacion: tambien define la regla local que implementaremos en el computador.
\end{alertblock}

\column{0.38\textwidth}
\begin{tikzpicture}[scale=0.85]
\foreach \x/\txt in {0/$u_{i-1}^n$,1.4/$u_i^n$,2.8/$u_{i+1}^n$} {
    \draw[fill=MidnightBlue!10] (\x,0) rectangle ++(1,1);
    \node at (\x+0.5,0.5) {\small \txt};
}
\node at (1.9,-0.6) {\Large $\Downarrow$};
\draw[fill=Crimson!12] (1.4,-2.0) rectangle ++(1,1);
\node at (1.9,-1.5) {\small $u_i^{n+1}$};
\end{tikzpicture}
\end{columns}
\end{frame}

\begin{frame}{Estabilidad numerica}
\justifying
El esquema explicito es simple y pedagogicamente muy util, pero no acepta cualquier eleccion de $\Delta t$.

\begin{block}{Resultado practico en 1D}
Para evitar inestabilidades, conviene exigir
\[
r=\frac{D\,\Delta t}{(\Delta x)^2} \le \frac{1}{2}.
\]
\end{block}

\begin{itemize}
    \item Si $r$ es demasiado grande, la simulacion puede oscilar y explotar numericamente.
    \item Si $r$ es muy pequeno, la simulacion sera estable pero puede ser lenta.
\end{itemize}
\end{frame}

\begin{frame}{Condiciones de borde}
\begin{itemize}
    \item \textbf{Dirichlet}: fijamos el valor de $u$ en los extremos.
    \item \textbf{Neumann}: fijamos el flujo, por ejemplo flujo nulo.
    \item \textbf{Periodicas}: el extremo izquierdo se conecta con el derecho.
\end{itemize}

\begin{block}{Eleccion pedagogica}
En esta semana es muy comodo usar bordes periodicos o flujo nulo, porque permiten concentrarse en la dinamica interna sin demasiadas complicaciones adicionales.
\end{block}
\end{frame}

\begin{frame}[fragile]{Codigo base: difusion 1D}
\inputminted[firstline=1,lastline=52]{python}{codes/diffusion_1d.py}
\end{frame}

\begin{frame}{Que esperar al correr el ejemplo 1D}
\begin{itemize}
    \item el perfil inicial se suaviza con el tiempo;
    \item el pico central disminuye;
    \item el ancho caracteristico aumenta;
    \item si el esquema es estable, la evolucion es suave y no aparecen oscilaciones espurias.
\end{itemize}
\end{frame}

\begin{frame}{De una dimension a dos dimensiones}
En 2D, la ecuacion es
\[
\pdv{u}{t} = D\left(\pdv[2]{u}{x} + \pdv[2]{u}{y}\right).
\]

El laplaciano discreto de cinco puntos es
\[
\nabla^2 u_{ij} \approx
\frac{u_{i+1,j}+u_{i-1,j}+u_{i,j+1}+u_{i,j-1}-4u_{ij}}{(\Delta x)^2}.
\]

\begin{block}{Comentario}
El paso conceptual es muy parecido al caso 1D: cambia la geometria del vecindario, no la idea central.
\end{block}
\end{frame}

\begin{frame}[fragile]{Codigo base: difusion 2D}
\inputminted[firstline=1,lastline=44]{python}{codes/diffusion_2d.py}
\end{frame}

\begin{frame}{Actividades propuestas para la Clase 1}
\begin{enumerate}
    \item Simular un pulso gaussiano 1D y comparar perfiles en distintos tiempos.
    \item Verificar empiricamente que un paso temporal demasiado grande vuelve inestable el esquema.
    \item Simular una mancha inicial 2D y observar su expansion.
    \item Medir una magnitud simple, por ejemplo el maximo de $u$ o la varianza espacial, como funcion del tiempo.
\end{enumerate}

\begin{block}{Entregable sugerido}
Script o notebook con el codigo, dos o tres graficos bien rotulados y una breve discusion sobre estabilidad y suavizado.
\end{block}
\end{frame}

\begin{frame}{Cierre Clase 1}
\begin{itemize}
    \item La difusion es un ejemplo directo de como una ley continua se vuelve una regla local en una malla.
    \item La estabilidad numerica no es un detalle tecnico: determina si la simulacion representa o no al sistema fisico.
    \item El paso natural siguiente es agregar reaccion local, es decir, produccion y destruccion de especies.
\end{itemize}
\end{frame}

\section{Clase 2}

\begin{frame}[plain]
\centering
\vfill
{\Huge \bfseries Clase 2}\\[0.4cm]
{\Large Reaccion-difusion, inestabilidad de Turing y patrones}
\vfill
\end{frame}

\begin{frame}{Objetivos de la Clase 2}
\begin{itemize}
    \item Entender que la combinacion de reaccion local y transporte puede producir patrones espaciales estables.
    \item Presentar la idea de inestabilidad de Turing sin sobrecargar el formalismo.
    \item Implementar una version discreta del modelo de Gray--Scott.
    \item Explorar parametros y medir rasgos cuantitativos simples de los patrones generados.
\end{itemize}
\end{frame}

\begin{frame}{De difusion a reaccion-difusion}
\justifying
Si ahora consideramos dos campos, por ejemplo $u(\bm{x},t)$ y $v(\bm{x},t)$, podemos combinar:

\begin{itemize}
    \item \textbf{difusion}: mezcla espacial;
    \item \textbf{reaccion}: conversion local entre especies;
    \item \textbf{alimentacion o perdida}: intercambio con el entorno.
\end{itemize}

La estructura general es
\[
\pdv{u}{t} = D_u \nabla^2 u + f(u,v),
\qquad
\pdv{v}{t} = D_v \nabla^2 v + g(u,v).
\]
\end{frame}

\begin{frame}{Idea de Turing}
\justifying
Alan Turing mostro que un estado uniforme puede ser estable frente a perturbaciones locales cuando no hay espacio, pero volverse inestable al combinarse con difusion.

\begin{block}{Interpretacion fisica}
La difusion no siempre destruye estructura. Si dos especies difunden a ritmos distintos y reaccionan entre si, puede amplificar ciertas escalas espaciales y producir bandas, manchas o laberintos.
\end{block}

\begin{itemize}
    \item no es necesario entrar en toda la teoria lineal para usarlo pedagogicamente;
    \item basta con entender que reaccion + transporte puede generar patron.
\end{itemize}
\end{frame}

\begin{frame}{Modelo de Gray--Scott}
Una forma muy usada en docencia y simulacion visual es
\begin{align*}
\pdv{u}{t} &= D_u \nabla^2 u - uv^2 + F(1-u), \\
\pdv{v}{t} &= D_v \nabla^2 v + uv^2 - (F+k)v.
\end{align*}

\begin{itemize}
    \item $u$ se alimenta desde el entorno;
    \item $v$ se consume y decae;
    \item el termino $uv^2$ acopla no linealmente ambas especies.
\end{itemize}
\end{frame}

\begin{frame}{Lectura cualitativa de los terminos}
\begin{columns}[T,onlytextwidth]
\column{0.58\textwidth}
\begin{itemize}
    \item $D_u \nabla^2 u$, $D_v \nabla^2 v$: suavizan gradientes.
    \item $-uv^2$: consume $u$ donde hay suficiente $v$.
    \item $+uv^2$: produce $v$ de forma autocatalitica.
    \item $F(1-u)$: repone $u$.
    \item $-(F+k)v$: retira $v$ del sistema.
\end{itemize}

\begin{alertblock}{Mensaje}
El patron no viene dibujado en la ecuacion. Aparece por la competencia entre mezcla, produccion y disipacion.
\end{alertblock}

\column{0.36\textwidth}
\begin{tikzpicture}[scale=0.92]
\node[draw,rounded corners,fill=MidnightBlue!10,minimum width=2.6cm,minimum height=0.9cm] (u) at (0,1.3) {$u$};
\node[draw,rounded corners,fill=Crimson!10,minimum width=2.6cm,minimum height=0.9cm] (v) at (0,-1.0) {$v$};
\draw[->,thick] (-2.2,1.3) -- (u) node[midway,above] {\small feed};
\draw[->,thick] (u) -- node[right] {\small $uv^2$} (v);
\draw[->,thick] (2.0,-1.0) -- (v) node[midway,above] {\small salida};
\end{tikzpicture}
\end{columns}
\end{frame}

\begin{frame}{Discretizacion computacional}
\justifying
La implementacion mas directa usa:
\begin{itemize}
    \item una malla 2D;
    \item un laplaciano discreto;
    \item actualizacion explicita para $u$ y $v$;
    \item una perturbacion local inicial para romper la uniformidad.
\end{itemize}

\begin{block}{Buena practica}
Separar el calculo del laplaciano y la actualizacion facilita luego cambiar parametros, condiciones de borde o el modelo de reaccion.
\end{block}
\end{frame}

\begin{frame}[fragile]{Codigo base: Gray--Scott en 2D}
\inputminted[firstline=1,lastline=77]{python}{codes/gray_scott.py}
\end{frame}

\begin{frame}{Exploracion de parametros}
\begin{itemize}
    \item Cambios pequenos en $(F,k)$ pueden cambiar mucho la morfologia.
    \item Algunos parametros producen manchas aisladas.
    \item Otros producen bandas, laberintos o patrones mas granulados.
    \item La resolucion espacial y el tiempo total de integracion tambien importan.
\end{itemize}

\begin{block}{Actividad natural}
Construir una pequena tabla de experimentos y describir visualmente el patron obtenido en cada caso.
\end{block}
\end{frame}

\begin{frame}{Mas alla de mirar imagenes}
\justifying
Un curso de modelacion computacional gana mucho cuando no se queda solo en la visualizacion. Tambien conviene medir.

\begin{itemize}
    \item fraccion del dominio donde $v$ supera cierto umbral;
    \item promedio y desviacion estandar de $u$ o $v$;
    \item longitud caracteristica estimada desde un perfil o una autocorrelacion simple;
    \item evolucion temporal de la varianza espacial.
\end{itemize}
\end{frame}

\begin{frame}[fragile]{Codigo de apoyo: metricas simples del patron}
\inputminted[firstline=1,lastline=30]{python}{codes/pattern_metrics.py}
\end{frame}

\begin{frame}{Actividades propuestas para la Clase 2}
\begin{enumerate}
    \item Simular Gray--Scott para un conjunto pequeno de pares $(F,k)$.
    \item Comparar visualmente al menos dos patrones distintos.
    \item Medir la fraccion activa y la desviacion estandar de $v$ en cada caso.
    \item Discutir si el patron parece converger a una estructura estable, oscilatoria o transitoria.
    \item Desafio: probar una condicion inicial distinta y comparar el resultado.
\end{enumerate}

\begin{block}{Entregable sugerido}
Script o notebook con imagenes, tabla corta de parametros y una interpretacion fisica/computacional del patron observado.
\end{block}
\end{frame}

\begin{frame}{Conexiones con el resto del curso}
\begin{itemize}
    \item Como en Ising y autómatas celulares, aparecen patrones globales desde reglas locales.
    \item Como en dinamica molecular y Monte Carlo, el algoritmo y la discretizacion afectan el resultado.
    \item Como antes del bloque de IA, esta semana deja un conjunto de datos espaciales ideal para futuras tareas de clasificacion, reduccion o aprendizaje de patrones.
\end{itemize}
\end{frame}

\begin{frame}{Cierre de la semana}
\begin{itemize}
    \item La reaccion-difusion muestra una forma muy rica de emergencia en sistemas continuos.
    \item Una discretizacion simple basta para producir experimentos computacionales muy expresivos.
    \item La Semana 08 cierra bien la unidad al combinar EDPs, simulacion 2D, estabilidad, visualizacion y analisis cuantitativo.
\end{itemize}
\end{frame}

\end{document}
